% Source: 06 - Wed Oct 7 2026.pdf, page 10. Content remains in source order.
\sourcepage{06}{10}
We now introduce a problem in the realm of COMPUTER ARITHMETIC:

\textbf{SUBSET-SUM (SS)}

I: $\langle S,t\rangle:S\subseteq\mathbb N$ finite, $t\in\mathbb N$ (``target'').

Q: $\exists S'\subseteq S:\sum_{s\in S'}s=t$?

The problem asks whether there is a subset of the input set whose sum is equal to $t$ (the target).

\textbf{EXAMPLE:}
\[\langle\{10,15,4,3,17\},24\rangle\in L_{\mathrm{SS}}\]
since $24=4+3+17$ ($S'=\{3,4,17\}$).

\textbf{REMARK:} $|\langle S,t\rangle|$ is the total number of bits needed to encode $s\in S\ \forall s$ and $t$:
\[|\langle S,t\rangle|=\Theta\Bigl(\sum_{s\in S}\log s+\log t\Bigr).\]
Clearly, $\mathrm{SS}\in\classNP$.

Candidate certificate: $y=\langle S'\rangle$, $S'\subseteq S$.

Let us prove that $\mathrm{SS}\in\classNPH$ by proving that $\mathrm{3\text{-}CNF\text{-}SAT}\polyred\mathrm{SS}$.
